\section{The substitutability spectrum in full}\label{si:spectrum}

This section records the CES results summarized in main-text Section~\ref{main-sec:ladder} and Theorem~\ref{main-thm:ladder}. Results concern the witness datum and the explicit collapse convention; $\sigma$ denotes the elasticity of substitution, while the Schaefer surplus in SI~S5 is written $\varsigma$.

\subsection{The critical-elasticity landscape on the diagonal}

For each gap state the accepted rungs form a prefix of the ladder (main-text Theorem~4(i)), so the $\sigma$-aggregator falsely certifies $z$ exactly when $\sigma\ge\sigma^*(z)$. Table~\ref{tab:sigmastar} lists exact brackets on the diagonal $x=1/2$; Figure~\ref{main-fig:ladder}(b) of the main text overlays them on the master-equation root.

\begin{table}[!htbp]
\centering\small
\caption{The critical elasticity $\sigma^*(z)$ on diagonal gap states $z=(1/2,s_1,s_2)$ (exact brackets; strict at both ends unless noted).}
\label{tab:sigmastar}
\begin{tabular}{@{}llp{6.4cm}@{}}
\toprule
state $z=(x,s_1,s_2)$ & $\sigma^*(z)$ & note \\
\midrule
$(1/2,\,1,\,1)$ & $\infty$ exactly & only the linear member certifies \\
$(1/2,\,3/2,\,1/2)$ & $\infty$ exactly & min coordinate $\le1$: only-linear \\
$(1/2,\,6/5,\,6/5)$ & $\in(2,3)$, strict & the witness point \\
$(1/2,\,5/4,\,5/4)$ & $=2$ exactly & the $\theta=1/2$ rung's equality state \\
$(1/2,\,13/10,\,13/10)$ & $\in(1,2)$ & \\
$(1/2,\,3/2,\,3/2)$ & $\in(1/2,1)$, strict & the LPI-form witness \\
$(1/2,\,13/8,\,13/8)$ & $\in(1/3,1/2)$ & the harmonic member's false-certification witness \\
$(1/2,\,9/5,\,9/5)$ & $\in(1/5,1/4)$ & \\
$(1/2,\,39/20,\,39/20)$ & $\in(1/17,1/13)$ & near the typed boundary: $\sigma^*\to0$ \\
\bottomrule
\end{tabular}
\end{table}

Gap states with $\min(s_1,s_2)\le1$ are only-linear ($\sigma^*=\infty$); states with both coordinates $>1$ have finite $\sigma^*$, given on the diagonal by the master-equation root, with $\sigma^*(s)\to0$ as $s\to2^-$ and $\sigma^*(s)\to\infty$ as $s\to1^+$. Concretely, $\sigma=3$ certifies $s=6/5$, $\sigma=2$ certifies $s=13/10$, $\sigma=1$ certifies $s=3/2$, $\sigma=1/2$ certifies $s=13/8$ and $\sigma=1/4$ certifies $s=9/5$ (the per-rung false-certification witnesses of main-text Theorem~4(iii)); the harmonic member's Fibonacci pair is $8/5$ rejects, $13/8$ accepts.

\subsection{The exact off-diagonal criterion}

On the viable gap quadrant $x<1$, $1<s_1,s_2<2$, the principal-plan trough inputs to the geometric member are positive. Write
\[
A=\ln(s_1-1),\quad B=\ln(s_2+1),\quad C=\ln(s_1+1),\quad D=\ln(s_2-1).
\]
The fast plan is licensed on a downward interval of weight share $p=w_1$, and the slow plan on an upward interval. Their union covers $[0,1]$ exactly when the interval endpoints overlap.

\begin{theorem}[off-diagonal geometric criterion]\label{thm:offdiag}
For $x<1$ and $1<s_1,s_2<2$, Protocol~2 at $\theta=0$ accepts exactly when
\[
AD\le BC,\qquad\text{equivalently}\qquad r(s_1)r(s_2)\le1,
\quad r(s)=\frac{\ln(1/(s-1))}{\ln(s+1)}.
\]
Equivalently, its acceptance region is $s_2\ge\rho(s_1)$, where $\rho=r^{-1}\circ(1/r)$ is a strictly decreasing involution. Its unique fixed point is $\sqrt2$, since $r(s)\le1$ iff $s^2\ge2$; its endpoint limits are $\rho(s)\to2^-$ as $s\to1^+$ and $\rho(s)\to1^+$ as $s\to2^-$. Thus $\min(s_1,s_2)\ge\sqrt2$ suffices for acceptance, while acceptance requires $\max(s_1,s_2)\ge\sqrt2$. For $\theta\in(0,1)$, define $A=(s_1-1)^\theta$, $B=(s_2+1)^\theta$, $C=(s_1+1)^\theta$, and $D=(s_2-1)^\theta$; the analogous interval-cover condition is exactly
\[
F:=C(B-1)+A(1-D)-(B-D)\ge0.
\]
These are pointwise criteria; they do not assert a global fixed-sum phase structure.
\end{theorem}

\begin{proof}
At $\theta=0$, the fast condition is $pA+(1-p)B\ge0$, hence $p\le U:=B/(B-A)$; the slow condition is $pC+(1-p)D\ge0$, hence $p\ge L:=-D/(C-D)$. All denominators are positive on the stated quadrant. The cover is $L\le U$, equivalent after cross-multiplication to $AD\le BC$. Since $A,D<0$ and $B,C>0$, this is $r(s_1)r(s_2)\le1$. The function $r$ is positive and strictly decreasing on $(1,2)$, so $\rho$ is well-defined and decreasing; the product identity gives $\rho(\rho(s))=s$. The fixed point is determined by $r(s)=1$, which is equivalent to $s^2=2$, and the limits follow from $r(s)\to\infty$ at $1^+$ and $r(s)\to0$ at $2^-$. The pivot statements follow because $r(s)\le1$ exactly for $s\ge\sqrt2$. For $0<\theta<1$, the same interval calculation gives $U=(B-1)/(B-A)$ and $L=(1-D)/(C-D)$; multiplying by the positive denominators and collecting terms yields $F\ge0$.
\end{proof}

\paragraph{An exact harmonic specialization} At $\theta=-1$, the interval-cover expression simplifies to
\[
F=\frac{2(s_1^2-s_1+s_2^2-s_2-2)}{(s_1-1)(s_1+1)(s_2-1)(s_2+1)}.
\]
The denominator is positive on $(1,2)^2$, so acceptance is equivalent to $(2s_1-1)^2+(2s_2-1)^2\ge10$. On the diagonal this reduces to $s\ge\varphi$, the harmonic critical floor. This identity follows algebraically.

\subsection{Scope of the off-diagonal results}

The criteria above give pointwise acceptance conditions. They do not establish the minimum or stationary-point structure on a fixed-sum slice, equality of off-diagonal minima, monotonicity of the critical elasticity, or a global phase diagram. No such global claim is used in the main text. The diagonal thresholds are proved analytically in SI~S2.

\subsection{The Living Planet Index identification, with its caveats}

The $\theta=0$ member has the geometric-mean functional form used by the Living Planet Index. Here it is applied to the constructed inputs $\lambda_i=1+s_i$, not to the published LPI abundance series. Each $s_i$ is normalized by its own declared positive reference scale; consequently the critical thresholds depend on that normalization. Under the optional biomass map in main-text Section~\ref{main-sec:benchmark}, $B_m/B_{\mathrm{lim}}=1+s_1/2$, whereas the constructed index is $\lambda_1=1+s_1$. They are not the same abundance ratio. The published LPI's input construction and fixed weighting are not Protocol~2; acceptance for every weight is more restrictive than acceptance at that one published weight, so the protocol's false-acceptance exposure can be smaller. No claim is made that the published LPI implements this assessment rule.

\subsection{Interpretation at the benchmark point}

At the heatwave trough, the constructed values are $\lambda=(1/5,11/5)$, giving an equal-weight linear mean $6/5$ and geometric mean $\sqrt{11/25}<1$. This is a normalized-margin illustration of how the aggregator can alter a report; it is not a claim about an operational LPI value or an automatic management response. The underlying biomass breach, terminal treatment, and income-margin assumptions are stated in SI~S5.

\section{The common-shock variant in full}\label{si:commonshock}

\subsection{The variant}

Decompose the witness datum per the benchmark arithmetic: each principal plan's announced schedule is a mid-interval tent drawdown of depth $3/2$ in its characteristic coordinate (\act{fast}: the quota pulse on $s_1$; \act{slow}: the gradual burden on $s_2$), the other coordinate constant; each environmental event is a transient tent dip of depth $\delta$ in one floor coordinate, troughing at $t=1/2$ and recovered by $t=1$, successors unchanged; the worst-case depth $2=3/2+1/2$ is the quota pulse plus the heatwave, so the datum's event depth is $\delta=1/2$. \act{staged} spends $x$ (it needs $x\ge1$) while its floors grow linearly to $s+e$; an event of depth $\delta$ on a growing floor dips it to $s+1/8-\delta$ -- by $\max(0,\delta-1/8)$. The \emph{common} class: the two coordinate events are separate disturbances, each striking every plan, so each plan's per-coordinate worst case takes its own schedule plus the event on that coordinate. The disturbance quantifier remains innermost. For \act{staged}, the aggregate dip $\max(0,\mathrm{raw}(p)-1/8)$ is piecewise linear in the weight share $p$. Admissibility is therefore the conjunction of two affine inequalities on the branches; since the other plans' dips are affine in $p$, each cover condition reduces to endpoint evaluations.

\paragraph{Relation to the action-indexed case} At $\delta=1/2$, the action-indexed classification and witness-point thresholds are those stated in main-text Theorem~2. The common- and coupled-shock variants below change the disturbance structure while retaining the declared plan menu and path semantics.

\subsection{The middle-cell verdict, the rescue trichotomy and blend fragility}

These are main-text Proposition~6(i)--(iii), proved in S2.4. Two remarks complete them. The gap is strictly smaller than $I$ by the exact $1/2$-cut: every action-indexed gap state with $\min(s_1,s_2)<1/2$ is weak-rejected under the common class (witness $(\tfrac12,\tfrac1{10},\tfrac{19}{10})$, action-indexed weak-accepted by $s_1+s_2=2$). And the blend window is class-conditional in general: under the action-indexed class at event depth $\delta$ it is nonempty exactly on $\{s_1+s_2\ge3/2+\delta\}$ (main-text Theorem~3 is the $\delta=1/2$ instance), under the common class exactly on $\{s_1\ge\delta,\ s_2\ge\delta,\ s_1+s_2\ge3/2+2\delta\}$; the coupled gap below, by contrast, is blend-closed.

\subsection{The coupled regime}

\begin{theorem}[the coupled regime, additive reading]\label{thm:coupled}
Under the coupled class -- one disturbance carrying both coordinate events simultaneously, riding on the plans' own schedules -- at event depth $\delta$: the typed accepted set coincides with the common class's at every depth, and
\begin{align*}
\Vweak^{\mathrm{cpl}}={}&\{x<1,\ s_1\ge\delta,\ s_2\ge\delta,\ s_1+s_2\ge\tfrac32+2\delta\}\\
&\cup\{x\ge1,\ \min(s_1,s_2)\ge\max(0,\delta-\tfrac18)\}.
\end{align*}
At the datum's full magnitude ($\delta=2$) every $I$-state is weak-rejected, but the claim is depth-dependent with exact boundary $\delta=5/4$: at heatwave magnitude ($\delta=1/2$) deep-$I$ states survive as coupled gap states (witness $(\tfrac12,\tfrac{19}{10},\tfrac{19}{10})$); no $I$-state survives at $\delta\ge5/4$, while $(\tfrac12,\tfrac{79}{40},\tfrac{79}{40})$ survives at $\delta=5/4-1/40$; and at full magnitude the separation relocates rather than vanishes: the coupled gap $\{x<1,\ 2\le s_1<\tfrac72,\ 2\le s_2<\tfrac72,\ s_1+s_2\ge\tfrac{11}2\}$ is nonempty (witness $(\tfrac12,3,3)$ is weak-accepted and typed-rejected).
\end{theorem}
\begin{proof}[Proof sketch]
Under the coupled class each plan's aggregate dip is the sum of both events' weighted contributions ($3p/2+\delta$ for \act{fast}) against the common class's maximum of single-event contributions. For $x<1$, the endpoint weights $p=0$ and $p=1$ impose the separate coordinate cuts $s_2\ge\delta$ and $s_1\ge\delta$; the interior kink imposes $s_1+s_2\ge3/2+2\delta$. The three conditions must be used together. For example, at $\delta=1/4$, $(s_1,s_2)=(1/10,39/20)$ passes the sum inequality but fails the $s_1\ge\delta$ cut. An $I$-state has $s_1,s_2<2$, hence sum $<4$, so no $I$-state survives when $\delta\ge5/4$; the states $(79/40,79/40)$ at $\delta=5/4-1/40$ satisfy both coordinate cuts and the sum bound. The witness states realize the stated boundaries.
\end{proof}

\begin{theorem}[the replace-reading collapse]\label{thm:replace}
Under the coupled-total class -- the replace reading, each announced-schedule plan's worst case being the full-magnitude dip $2$ in each floor, tent schedules suppressed -- \act{fast} and \act{slow} degenerate to the same weight-independent requirement $A(p)\ge2$ (at the witness point both serve nowhere; at $(\tfrac12,\tfrac52,\tfrac52)$ both serve all of $[0,1]$; at $(\tfrac12,3,1)$ both serve exactly $[\tfrac12,1]$), and
\[
\Vweak^{\mathrm{ct}}=\Vtyp^{\mathrm{ct}}=\{s_1\ge2,\ s_2\ge2\}\cup\{x\ge1,\ \min(s_1,s_2)\ge\tfrac{15}8\}:
\]
the typed/weak separation collapses entirely, and the replace reading is strictly more demanding than the action-indexed class (the witness $(\tfrac12,3,1)$ is action-indexed typed-accepted yet a full reject here, as is the witness point itself).
\end{theorem}
\begin{proof}[Proof sketch]
With the tents suppressed and the full dip $2$ in each floor, every principal plan's aggregate dip is the constant $2$, so the weak cover condition $A(p)\ge2$ for all $p$ is the pair of endpoint conditions $s_1\ge2$, $s_2\ge2$, which is the typed condition; \act{staged}'s branch (the dip lands on the growing floor, threshold $15/8$) is identical under both tests.
\end{proof}

The two readings are distinct stylized disturbance specifications; neither is established as an upper or lower bound on a physically realized coupled-shock process. They illustrate model sensitivity, not a bracket on empirical truth. Under the replace reading, the plan- and weight-independent dip removes the weak test's advantage of weight-wise plan choice, so the per-weight cover reduces to the per-coordinate condition used by the typed test.

\paragraph{Zero-disturbance case} At $\delta=0$, the disturbance classes coincide in this stylized datum.

\paragraph{Model-relative consequences at the benchmark datum} These are conditional outputs of the constructed menu, disturbance classes, and path semantics, not empirical estimates or management prescriptions. The common-shock rule cuts the action-indexed weak gap at $\min(s)=1/2$ (for example, $(\tfrac12,\tfrac1{10},\tfrac{19}{10})$); no reserve increment rescues failure states below the $3/8$ floor cut. The coupled additive and replace readings yield distinct verdicts, neither claimed to bound a physical process. Under the common-shock rule, blended quota policies also fail on the specified band $s_1+s_2\in[2,\tfrac52)$. These are menu-relative feasibility results, not operational advice.

\paragraph{Verification and limits} The common- and coupled-class classifications follow from the analytic derivations above and in SI~S2.4. Their scope is the specified menu, disturbance classes, and path semantics; they are not empirical estimates or bounds on realized shock processes.

\begin{table}[!htbp]
\centering\small
\caption{Disturbance regimes on the witness datum and the status of the separation under each.}
\label{tab:disturbance}
\begin{tabularx}{\textwidth}{@{}>{\raggedright\arraybackslash}p{0.2\textwidth}>{\raggedright\arraybackslash}p{0.3\textwidth}X@{}}
\toprule
Regime & Disturbance class & Separation status \\
\midrule
(i) Action-indexed separable & each plan's worst case is its own characteristic dip in the coordinate it moves & proven: main-text Theorem~2 \\
(ii) Common shock & one shared shock strikes the same floor coordinate of every plan & computed (main-text Proposition~6): the gap survives, reduced by the exact $1/2$-cut, licensing thresholds unshifted on the surviving region; the witness point remains a member \\
(iii) Coupled all-floor shocks & simultaneous dips across all floors of every plan & reading-conditional (Theorems~\ref{thm:coupled}--\ref{thm:replace}): additive -- universal rejection exact on $I$ at full magnitude, depth-dependent below $\delta=5/4$, the gap relocating at full magnitude; replace -- the typed/weak distinction itself collapses, with acceptance surviving on $\{s_1\ge2,s_2\ge2\}\cup\{x\ge1,\min(s)\ge15/8\}$ \\
\bottomrule
\end{tabularx}
\end{table}
